\documentclass{article}
\usepackage[UTF8]{ctex}
\usepackage{amsmath,amssymb,amsfonts}
\usepackage{enumitem}
\usepackage{geometry}
\geometry{a4paper, margin=1in}
% === 错题本排版优化（不参与题目标识） ===
\setlength{\parindent}{0pt}
\setlength{\parskip}{0.35em plus 0.10em minus 0.08em}
\setlength{\emergencystretch}{3em}
\setlist{itemsep=0.25em, topsep=0.40em, parsep=0pt}
\allowdisplaybreaks[2]
% === 错题本排版优化结束 ===
\begin{document}
\section*{华东师范大学2024年数学分析真题}
\section*{36．华东师范大学2024年数学分析试卷}

\subsection{题目1}
1．（10 分）计算曲面 \(z=e^{x^{2}-y^{2}}+x^{2}\) 与曲面 \(z=e^{x^{2}-y^{2}}-y^{2}+1\) 所围成的封闭区域的体积．

\subsection{题目2}
2．（10 分）设 \(f(x)\) 是 \(\mathbb{R}\) 上以 \(2 \pi\) 为周期的连续函数，以 \(a_{n}, b_{n}\) 为其傅里叶级数的系数，求

\[
F(x)=\frac{1}{\pi} \int_{-\pi}^{\pi} f(t) f(x+t) \mathrm{d} t
\]

的傅里叶系数（用 \(a_{n}, b_{n}\) 表示）．

\subsection{题目3}
3．（15 分）设 \(\lim _{n \rightarrow \infty} \frac{a_{1}+a_{2}+\cdots+a_{n}}{n}=a(-\infty<a<+\infty)\) ．\\
（1）判断 \(\lim _{n \rightarrow \infty} a_{n}\) 是否存在，若存在，请给出严格证明；若不存在，请举例并详细说明．\\
（2）若 \(\lim _{n \rightarrow \infty} n\left(a_{n}-a_{n-1}\right)=0\) 或 \(\left\{a_{n}\right\}\) 单调递增，判断 \(\lim _{n \rightarrow \infty} a_{n}\) 是否存在，若存在，请给出严格证明；若不存在，请举例并给出详细证明．

\subsection{题目4}
4．（15 分）设 \(f(x)\) 在 \([0,1]\) 上有连续的导数，\(f(0)=0\) 或 \(\int_{0}^{1} f(x) \mathrm{d} x=0\) ．判断不等式

\[
\int_{0}^{1}[f(x)]^{2} \mathrm{~d} x \leq \int_{0}^{1}\left[f^{\prime}(x)\right]^{2} \mathrm{~d} x
\]

是否成立，若成立，请给出严格证明；若不成立，请举例并给出详细证明．

\subsection{题目5}
5．（15 分）讨论 \(\sum_{n=1}^{\infty} \frac{(-1)^{n-1}}{\left[n+(-1)^{n-1}\right]^{p}}(p>0)\) 的敛散性（收敛时需指明是绝对收敛还是条件收敛）。

\subsection{题目6}
6．（20 分）解答如下问题：

\begin{itemize}
  \item[（1）] 证明：\(\left\{u_{n}(x)\right\}=\left\{\left(1+\frac{x}{n}\right)^{n}\right\}\) 在 \(x \in[0,1]\) 上一致收敛。
  \item[（2）] 证明：\(\left\{v_{n}(x)\right\}=\left\{\frac{1}{e^{\frac{x}{n}}+\left(1+\frac{x}{n}\right)^{n}}\right\}\) 在 \(x \in[0,1]\) 上一致收敛。
  \item[（3）] 求 \(\lim _{n \rightarrow \infty} \int_{0}^{1} \frac{\mathrm{~d} x}{e^{\frac{x}{n}}+\left(1+\frac{x}{n}\right)^{n}}\) ．
\end{itemize}

\subsection{题目7}
7．（25 分）设 \(f\) 是定义在 \((-\infty,+\infty)\) 上的非常数函数，请判断下列说法是否正确．若正确，请给出严格证明；若不正确，请举例并给出详细说明．

\begin{itemize}
  \item[（1）] 若 \(f\) 是处处不连续的周期函数，则 \(f\) 必有最小正周期．
  \item[（2）] 若 \(f\) 是处处不连续的周期函数，则 \(f\) 必没有最小正周期．
  \item[（3）] 若 \(f\) 是周期函数但没有最小正周期，则 \(f\) 必有一列趋于 0 的周期．
  \item[（4）] 若 \(f\) 是连续的周期函数，则 \(f\) 必有最小正周期．
\end{itemize}

\subsection{题目8}
8．（20 分）设 \(f\) 在 \((a,+\infty)\) 上有连续的导数且 \(\lim _{x \rightarrow+\infty} f(x)=A(-\infty<A<+\infty)\) ．

\begin{itemize}
  \item[（1）] 请给出例子详细说明 \(\lim _{x \rightarrow+\infty} f^{\prime}(x)\) 末必存在。
  \item[（2）] 若 \(\lim _{x \rightarrow+\infty} f^{\prime}(x)\) 存在，求它的值．
  \item[（3）] 若 \(f(x)\) 满足 \(n\) 阶可导且 \(\lim _{x \rightarrow+\infty} f(x)\) 与 \(\lim _{x \rightarrow+\infty} f^{(n)}(x)\) 存在．判断对任意的 \(k=1,2, \cdots, n-1\) ，\(\lim _{x \rightarrow+\infty} f^{(k)}(x)\) 是否存在．若存在，请求出极限；若不存在，请举例并给出详细说明．
\end{itemize}

\subsection{题目9}
9．（20 分）设 \(f: \mathbb{R}^{2} \rightarrow \mathbb{R}\) 连续可微，\(f(0,0)=0\) 且 \(f_{y}(0,0) \neq 0\) ．考虑迭代：

\[
y_{n+1}(x)=y_{n}(x)-\left(f_{y}(0,0)\right)^{-1} f\left(x, y_{n}(x)\right), y_{0}(x)=0 .
\]

定义 \(u_{n+1}(x)=y_{n+1}(x)-y_{n}(x), n \geq 0\) ．\\
（1）证明：对任何 \(\varepsilon>0\) ，存在 \(\delta>0\) ，当 \(x \in(-\delta, \delta)\) 时，\(\left|y_{n}(x)\right|<\varepsilon\) 对任何 \(n \geq 0\) 成立，进而当 \(\delta\) 充分小时，存在 \(0<\rho<1\) 使得对任意的 \(|x|<\delta\) ，有 \(\left|u_{n+1}(x)\right| \leq \rho\left|u_{n}(x)\right|\) ．\\
（2）证明：\(y(x)=\sum_{n=1}^{\infty} u_{n}(x)\) 时 \(f(x, y)=0\) 在 \((0,0)\) 附近确定的一个连续的隐函数．
\end{document}
